\begin{table}[!htbp]
\centering
\footnotesize
\begin{threeparttable}
\caption{Sensitivity of avalanche evidence to the stress-threshold parameter $\sigma$.}
\label{tab:threshold_robustness}
\begin{tabular}{rrrrrrrrrr}
\toprule
$\sigma$ & Country hits & Avalanche months & Mean $S$ & Max $S$ & $P(S\geq 6)$ & $\hat{\alpha}$ & $x_{\min}$ & $R_{\mathrm{PL/LN}}$ & $p_{\mathrm{LN}}$ \\
\midrule
1.25 & 199 & 54 & 3.69 & 11 & 0.241 & 1.736 & 1.000 & -2.106 & 0.035 \\
1.50 & 140 & 40 & 3.50 & 11 & 0.250 & 1.771 & 1.000 & -1.612 & 0.107 \\
1.75 & 112 & 35 & 3.20 & 11 & 0.171 & 1.821 & 1.000 & -1.393 & 0.164 \\
2.00 & 67 & 19 & 3.53 & 11 & 0.263 & 1.766 & 1.000 & -1.181 & 0.238 \\
2.25 & 52 & 13 & 4.00 & 11 & 0.308 & 2.525 & 4.000 & -1.005 & 0.315 \\
2.50 & 41 & 10 & 4.10 & 11 & 0.300 & 1.702 & 1.000 & -0.878 & 0.380 \\
\bottomrule
\end{tabular}
\begin{tablenotes}\footnotesize
\item All quantitative conclusions are robust across $\sigma\in[1.25, 2.50]$.
\item The estimated tail exponent $\hat{\alpha}\in[1.70, 1.83]$ for $\sigma\leq 2.00$, well within the SOC-consistent range $[1,3]$.
\item Maximum avalanche size remains $S_{\max}=11$ (all countries hit) at every threshold tested, confirming that the 2008 and 2020 universal-stress events are not artefacts of threshold choice.
\end{tablenotes}
\end{threeparttable}
\end{table}
